\chapter{Einleitung}
\label{ch:einleitung}

Herzlich willkommen zum Handbuch für das \emph{(be)greifbare Smart Home}. Dieses Dokument dient als praxisorientierter Leitfaden und ist als Begleitmaterial der vorliegenden Masterarbeit konzipiert. Es richtet sich an alle, die das vorgestellte lokale, sichere und vollständig transparente Smart-Home-System für den Schulunterricht von Grund auf selbst aufbauen möchten.

Im Gegensatz zu proprietären Lösungen großer Hersteller steht hier der Gedanke der digitalen Souveränität im Vordergrund. Die Technologie soll im wahrsten Sinne des Wortes „begreifbar“ und kontrollierbar bleiben. In den folgenden Kapiteln wird Schritt für Schritt durch den gesamten Einrichtungsprozess geführt. Dieser reicht von der initialen Vorbereitung der Hardware und dem Flashen des Betriebssystems über die Installation des Xubuntu-Servers bis hin zur sicheren Netzwerkkonfiguration (SSH) und der finalen Inbetriebnahme der Smart Home Dienste und Geräte.

Im Handbuch befindet sich auch ein Vorschlag für eine konkrete praktische Umsetzung als modulare Smart Home Rauminstallationen. Alle Schritte lassen sich aber auch unabhängig vom vorgeschlagenen Formfaktor durchführen.

Dank detaillierter Code-Listings, Terminal-Ausgaben und kommentierter Konfigurationsdateien ist dieses Handbuch so strukturiert, dass der Aufbau der Infrastruktur systematisch nachvollzogen und reproduziert werden kann. Alle Config-Skripts und Codebeispiele sind in einem eigenen GitHub-Repository verfügbar. Dort finden sich auch weiterführende Informationen, Updates und sowie die begleitende schriftliche Masterarbeit zum Thema.

\vspace{2em}

\begin{InfoBox}{Begleitendes GitHub-Repository}
\url{https://github.com/Johannes-Baischer/Das--be-greifbare-Smart-Home}
\end{InfoBox}

\vspace{2em}

\begin{InfoBox}{Wichtiger Hinweis zu Abbildungen und Quellcode}
\scriptsize
Soweit nicht in der jeweiligen Bild- oder Codeunterschrift explizit anders angegeben, handelt es sich bei allen Screenshots, Diagrammen und Terminal-Ausgaben in diesem Handbuch um eigene Darstellungen des Verfassers, die im Rahmen der praktischen Umsetzung dieser Masterarbeit erstellt wurden.
\end{InfoBox}